% Capítulo 1 — Introdução (orçamento: ~8 pp)
% Voz: primeira pessoa quando ajuda, frases curtas, uma ideia por frase.
% Cada secção abre com a pergunta a que responde.

\chapter{Introdução}
\label{cap:introducao}

\section{De onde veio este trabalho}
\label{sec:intro_origem}

\emph{Que problema é este, e porque é que vale a pena resolvê-lo?}

Hoje é fácil comprar ações. Basta uma aplicação no telemóvel e alguns minutos. A maior parte das
famílias americanas já tem ações, diretamente ou através de fundos e planos de reforma
\autocite{gallup2025stock}, e o mercado onde investem vale mais de 62 biliões de dólares
\autocite{sifma2025factbook} (Figura~\ref{fig:mercado_eua}).

Difícil é a parte seguinte: perceber o que aconteceu quando o preço se mexe.

Imagine-se uma pessoa que tem cinco ou seis ações e olha para elas de vez em quando. Um dia abre a
aplicação e vê uma delas a cair 4\%. A aplicação diz-lhe exatamente isso: menos 4\%. E mais nada.

A partir daí a pessoa fica sozinha com três perguntas, e faz sempre estas três, por esta ordem:

\begin{enumerate}
  \item \textbf{Isto é invulgar?} Quatro por cento numa empresa calma é muito. Na mesma semana, e
        noutra empresa, pode ser um dia perfeitamente normal.
  \item \textbf{Foi a empresa, ou foi o mercado?} Se caiu tudo, a história é uma. Se caiu só esta,
        é outra completamente diferente.
  \item \textbf{Já aconteceu antes, e o que se seguiu?} Se já houve dias parecidos, o que fez o
        preço a seguir?
\end{enumerate}

As aplicações gratuitas não respondem a nenhuma das três. Mostram a percentagem, que é o princípio
da primeira pergunta e não a resposta. Os terminais profissionais respondem às três, mas custam mais
por mês do que muita gente investe por ano.

É essa distância que este trabalho tenta encurtar. Não por falta de dados: há dados a mais, e
gratuitos. O que falta é \textbf{contexto}, e um contexto que a pessoa possa \emph{verificar}.

\begin{figure}[ht]
\centering
\includegraphics[width=0.85\textwidth]{figures/us_equity_market_cap.pdf}
\caption[Capitalização do mercado acionista dos EUA]{Capitalização do mercado acionista dos EUA,
2015--2024, em biliões de dólares. É neste mercado que investem as pessoas de quem este trabalho
fala. Elaborado a partir dos dados do SIFMA Capital Markets Fact Book
\autocite{sifma2025factbook}.}
\label{fig:mercado_eua}
\end{figure}

\section{O que o sistema faz, numa página}
\label{sec:intro_sistema}

\emph{Se tivesse de explicar o InvestiGator numa frase, qual seria?}

Esta: \textbf{o sistema vigia um conjunto de empresas, deteta quando algo sai fora do normal, e
avisa, mas só avisa com a explicação e a evidência ao lado.}

Há duas maneiras de o sistema começar a prestar atenção, e estão na
Figura~\ref{fig:conceito}. Ou o \textbf{preço} se mexe de forma invulgar, ou chega uma
\textbf{notícia} relevante sobre uma das empresas. Em qualquer dos casos, o que sai do outro lado
não é uma notificação seca. É um alerta que diz o que aconteceu, porque é que foi considerado
invulgar, de onde vieram os dados, e que casos parecidos já existiram no passado, com o que o
preço fez depois deles.

\begin{figure}[ht]
\centering
\begin{tikzpicture}[
  font=\small, >={Stealth[length=2.2mm]},
  gat/.style={draw, rounded corners, fill=black!8, align=center, text width=3.1cm, minimum height=1cm},
  sis/.style={draw, rounded corners, fill=black!3, align=center, text width=2.6cm, minimum height=1.6cm},
  saida/.style={draw, rounded corners, fill=black!3, align=center, text width=3.8cm, minimum height=1.6cm},
]
  \node[gat] (g1) {movimento invulgar\\ no preço};
  \node[gat, below=7mm of g1] (g2) {notícia relevante\\ sobre a empresa};
  \node[sis] (sis) at ($(g1)!0.5!(g2) + (4.2,0)$) {\textbf{InvestiGator}};
  \node[saida, right=1.3cm of sis] (out) {alerta explicado\\[2pt]
    \footnotesize o que aconteceu\\ \footnotesize porquê · fontes\\ \footnotesize casos parecidos};
  \node[right=1.0cm of out, align=center] (inv) {\footnotesize investidor};
  \draw[->] (g1) -- (sis);
  \draw[->] (g2) -- (sis);
  \draw[->] (sis) -- (out);
  \draw[->] (out) -- (inv);
\end{tikzpicture}
\caption[O que o sistema faz]{Os dois gatilhos e o que sai deles. O que o investidor recebe é um
alerta que ele consegue seguir e conferir, e não apenas um aviso de que o preço mexeu.}
\label{fig:conceito}
\end{figure}

\subsection*{E o que o sistema recusa fazer}

Há uma coisa que este sistema não faz, e não é por não ter dado jeito: \textbf{não prevê preços}.

Nunca diz que uma ação vai subir. Nunca diz para comprar ou para vender. Nunca mostra um alvo de
preço.

Isto pode parecer uma limitação. É o contrário: é o que torna tudo o resto verificável. Uma
afirmação sobre o que \emph{já} aconteceu pode ser confirmada contra o registo no momento em que é
lida. Uma previsão não pode: só se saberá se estava certa depois de o investidor já ter tido de
decidir. Ao recusar prever, o sistema obriga-se a dizer apenas aquilo que consegue mostrar.

\section{Objetivos e perguntas de investigação}
\label{sec:intro_objetivos}

\emph{O que é que este trabalho se propõe descobrir?}

O objetivo geral é \textbf{desenhar, construir e avaliar} um sistema de alertas financeiros
explicável para investidores de retalho, usando apenas fontes de dados gratuitas.

Para lá chegar, o trabalho responde a três perguntas de investigação. Cada uma tem uma medição
associada, porque uma pergunta a que não se sabe responder com um número não é uma pergunta de
investigação: é uma opinião.

\begin{itemize}
  \item \textbf{\gls{QI}1, deteção.} Consegue um método estatístico simples e transparente
        detetar movimentos invulgares de forma consistente entre empresas com volatilidades muito
        diferentes, e fazê-lo melhor do que uma regra fixa?
  \item \textbf{\gls{QI}2, precedentes.} Consegue o sistema encontrar notícias passadas
        genuinamente parecidas com a atual, melhor do que alternativas triviais, e medir o que
        aconteceu ao preço a seguir a elas sem usar informação do futuro?
  \item \textbf{\gls{QI}3, triagem.} Consegue um modelo treinado decidir que notícias merecem um
        alerta melhor do que uma linha de base simples baseada apenas em volatilidade?
\end{itemize}

Estas três perguntas dão origem a quatro objetivos concretos:

\begin{enumerate}
  \item construir o sistema completo, dos dados até ao alerta entregue;
  \item avaliar cada componente contra linhas de base transparentes;
  \item garantir que as explicações são fiéis ao que o sistema realmente calculou;
  \item reportar os resultados tal como se revelarem, incluindo os que forem negativos.
\end{enumerate}

O quarto objetivo é o mais importante e é o mais fácil de trair. Voltarei a ele no
Capítulo~\ref{cap:conclusoes}.

\section{Contribuições}
\label{sec:intro_contribuicoes}

\emph{O que é que fica feito no fim?}

Esta é uma dissertação de \textbf{Engenharia de Inteligência Artificial}. A contribuição não é
inventar um algoritmo novo. É escolher, integrar e avaliar criticamente técnicas que já existem, até
formarem um sistema que funciona, que se explica, e que outra pessoa consegue reproduzir. Escolher
bem \emph{é} o trabalho de engenharia.

Ficam quatro contribuições:

\begin{enumerate}
  \item \textbf{Um sistema completo e a funcionar}, com os dois gatilhos, sobre fontes gratuitas, e
        implantado, e não uma demonstração de laboratório.
  \item \textbf{Um método reprodutível para ligar notícias a impacto de mercado}, que combina
        recuperação por semelhança semântica com janelas de estudo de evento, sem deixar entrar
        informação do futuro.
  \item \textbf{Um conjunto de dados próprio e um modelo treinado sobre ele}: 79\,753 exemplos
        construídos a partir de notícias e preços históricos, com o modelo, a calibração e as
        métricas guardados como artefactos versionados.
  \item \textbf{Uma avaliação honesta de cada componente} contra linhas de base transparentes,
        incluindo o resultado que me correu ao contrário, que está reportado como tal.
\end{enumerate}

\section{Como ler este documento}
\label{sec:intro_estrutura}

\emph{Por onde é que isto vai?}

O documento foi escrito para ser lido do princípio ao fim, e cada capítulo responde a uma pergunta:

\begin{itemize}
  \item \textbf{Contexto} (Capítulo~\ref{cap:contexto}): o que já existe, e porque é que não
        chega para o problema desta tese.
  \item \textbf{Metodologia} (Capítulo~\ref{cap:metodologia}): as três técnicas, explicadas
        devagar e com desenhos.
  \item \textbf{O sistema} (Capítulo~\ref{cap:sistema}): como as peças se ligam, e uma notícia
        real seguida do princípio ao fim.
  \item \textbf{Avaliação} (Capítulo~\ref{cap:avaliacao}): o que foi medido, contra o quê, e com
        que resultado.
  \item \textbf{Conclusões} (Capítulo~\ref{cap:conclusoes}): o que ficou provado, o que ficou por
        fazer, e o que eu faria a seguir.
\end{itemize}

Há um fio que atravessa os capítulos~\ref{cap:metodologia} a~\ref{cap:avaliacao}: \textbf{a mesma
notícia real}, seguida desde o momento em que chega até ao alerta que sai. Sempre que aparecer, é a
mesma. Serve para que nada aqui fique abstrato.
